% ==================================================================
% Section 5 --- Results  (updated after DeepSeek zero-shot CSV recovery)
% Copy-paste this entire block into eisej_paper.tex, replacing the
% current Section 5.  All formatting is proper LaTeX (no markdown).
% All numbers are from the eval driver's final CSVs.
% ==================================================================

\section{Results}\label{sec:Results}

This section reports the evaluation of the four models across the
three research questions defined in Section~\ref{sec:Introduction}.
Section~\ref{sec:results-rq1} addresses RQ1 (acceptance
classification), Section~\ref{sec:results-rq2} addresses RQ2 (overall
flag-generation performance), and Section~\ref{sec:results-rq3}
addresses RQ3 (extractive versus evaluative flag categories).
Section~\ref{sec:results-drift} presents a supplementary finding on
vocabulary drift under few-shot prompting.

\subsection{RQ1: Acceptance Classification}\label{sec:results-rq1}

Table~\ref{tab:acceptance} reports acceptance-classification
performance across all four models and prompting conditions, together
with a majority-class baseline that always predicts \emph{merged} ---
the majority label in the dataset (206 merged, 94 unmerged; 68.7\%
accuracy floor).  Reporting this baseline separates genuine predictive
capability from the effect of class imbalance.

% >>> PASTE TABLE 5.1 HERE  (table_5_1_acceptance.tex) <<<

All four models exceed the majority baseline on overall accuracy,
though by markedly different margins.  \textbf{Claude Haiku 4.5}
achieves the highest single-condition accuracy (0.860 on
one-shot-unmerged), followed closely by its zero-shot performance
(0.853), representing a 16.6- and 16.5-point improvement over baseline
respectively.  Across all four prompting conditions, Haiku 4.5 remains
between 0.806 and 0.860, exhibiting the narrowest performance spread
of any model.  \textbf{DeepSeek R1} achieves consistently high
accuracy (0.801--0.825) with a distinctive profile: near-perfect
recall on merged PRs (0.99--1.00) at the cost of substantially lower
recall on unmerged PRs (0.37--0.46), indicating a systematic bias
toward predicting acceptance.  \textbf{GPT-4o-mini} shows the largest
variance across prompting conditions (0.716--0.774) and is the only
model whose few-shot condition (0.774) does not exceed its zero-shot
(0.773).  \textbf{Mistral Small 3.2-24B Instruct} falls in a similar
band (0.773--0.808), with one-shot-unmerged (0.799) and few-shot
(0.808) modestly improving on zero-shot.

Improvement over baseline is stronger on macro-averaged F1 than on
raw accuracy, since the baseline achieves 0.00 F1 on the minority
(unmerged) class.  Every non-baseline model--condition combination
attains F1(No)~$>$~0.5, indicating that all models learn some
discriminative capacity for unmerged PRs --- a signal the accuracy
metric alone would understate given the class imbalance.

Two rows of Table~\ref{tab:acceptance} report reduced effective~$N$:
DeepSeek R1 one-shot-merged ($N=297$, 2 excluded) and one-shot-unmerged
($N=298$, 1 excluded).  These excluded rows are cases where the model
returned no parseable JSON output despite retry, and are dropped from
evaluation per the exclusion policy documented in
Section~\ref{sec:Method}.

\subsection{RQ2: Flag Generation (Multi-Label)}\label{sec:results-rq2}

Flag generation is a substantially harder task than acceptance
prediction.  Table~\ref{tab:multilabel} reports the six multi-label
metrics defined in Section~\ref{sec:Method} across the same
$4 \times 4$ grid.

% >>> PASTE TABLE 5.2 HERE  (table_5_2_multilabel.tex) <<<

Across all sixteen model--condition combinations, macro-averaged F1
ranges from 0.003 (GPT-4o-mini few-shot) to 0.543 (DeepSeek R1
one-shot-unmerged) --- an order-of-magnitude spread that indicates
flag generation quality is not primarily a function of model quality,
but interacts strongly with prompting condition.  \textbf{DeepSeek R1
produces the strongest overall multi-label performance}, achieving the
highest macro-F1 (0.543 on one-shot-unmerged), highest micro-F1
(0.774 on one-shot-merged), and lowest Hamming loss (0.138 on
one-shot-merged) across the grid.  \textbf{Claude Haiku 4.5} exhibits
the most stable multi-label profile: macro-F1 varies only between
0.418 and 0.454 across all four conditions, and micro-F1 remains
between 0.734 and 0.762 --- indicating that in-context examples
neither meaningfully improve nor degrade its flag-generation
capability.

The two smallest-capacity models exhibit a pathological pattern under
few-shot prompting: \textbf{GPT-4o-mini's macro-F1 drops from 0.415
(zero-shot) to 0.003 (few-shot)}, and \textbf{Mistral's from 0.399 to
0.016}.  This is not a subtle degradation; it is a near-total collapse
of measurable flag-generation performance.
Section~\ref{sec:results-drift} examines the cause of this collapse.

Exact-match accuracy --- the strictest metric, requiring every flag on
a PR to be predicted correctly with no additions or omissions ---
remains below 0.06 for every model--condition combination.  This
reflects the intrinsic difficulty of multi-label prediction over 24
flags: even a model achieving 0.9 accuracy on each flag independently
would achieve $0.9^{24} \approx 0.08$ on exact match.  The metric is
included for completeness but should not be over-interpreted.

\subsection{RQ3: Extractive vs.\ Evaluative Flag
Categories}\label{sec:results-rq3}

To address RQ3, we partition the 24-flag vocabulary into two
categories defined by the prompt itself.  \textbf{Extractive flags}
(4 total) are those whose applicability the prompt states in terms of
an explicit numeric threshold applied to a field already present in
the prompt: \texttt{low\_contributor\_acceptance\_rate}
($<50\%$), \texttt{high\_contributor\_acceptance\_rate}
($\geq 80\%$), \texttt{large\_number\_of\_changed\_lines}
($>500$), and \texttt{manageable\_number\_of\_changed\_lines}
($<300$).  These flags require no semantic judgment --- a model that
reads the relevant number and correctly applies the stated threshold
will label them correctly.  \textbf{Evaluative flags} (the remaining
20) require the model to reason about semantic content --- code
quality, alignment with the linked issue, sentiment of review
comments, clarity of PR title or description --- with no explicit
thresholds supplied.

Table~\ref{tab:extractive-evaluative} reports the average per-flag F1
within each category, for every model--condition combination.

% >>> PASTE TABLE 5.3 HERE  (table_5_3_extractive_vs_evaluative.tex) <<<

The extractive-vs-evaluative gap is present across every model and
every non-degenerate prompting condition, ranging from 0.20 (Mistral
one-shot-unmerged) to 0.49 (GPT-4o-mini zero-shot).  Extractive flags
are identified reliably: DeepSeek R1 achieves 0.82--0.91 average F1
on extractive flags across the four conditions, Haiku 4.5 achieves
0.69--0.79, GPT-4o-mini 0.53--0.78 (excluding its few-shot collapse),
and Mistral 0.47--0.67.  In contrast, evaluative flags are identified
only weakly: no model exceeds 0.47 average F1 on evaluative flags in
any condition, and most fall between 0.30 and 0.44.  This gap is the
paper's central RQ3 finding.

Inspection of the per-flag detail (per-model tables in the online
supplementary material) reveals that within the evaluative category,
four flags --- \texttt{poor\_code\_quality},
\texttt{overengineered\_solution},
\texttt{solution\_incorrect\_or\_inefficient}, and to a lesser extent
\texttt{responsive\_to\_feedback} --- score below F1~$=0.10$ for
essentially every model and condition.  These four flags all concern
subjective judgments about code correctness, design appropriateness,
or contributor behavior across the review lifecycle: assessments a
human reviewer typically makes only after sustained engagement with
the code and the discussion around it, drawing on background context
that is not fully recoverable from the prompt alone.  That current
LLMs at this scale cannot reliably produce these judgments ---
regardless of the number of in-context examples provided --- is
consistent with the pattern observed by prior empirical work on PR
merge decisions~\cite{Zhang2023,Lenarduzzi2021} that social and
process-level signals often dominate code-quality signals in the
actual accept/reject determination.

Extractive flags, in contrast, correspond to structured criteria for
which the ground-truth is mechanically computable and directly present
in the prompt.  Their high F1 confirms that the models correctly parse
and apply the numeric thresholds provided, but is not itself evidence
of reasoning capacity --- a rule-based system would achieve F1~$=1.0$
on all four extractive flags without any language model.  The paper's
principal empirical claim is therefore that, for the
interpretable-flag-generation task explored here, LLM-based approaches
deliver reliable performance only on the extractive subset --- the
same subset most amenable to non-LLM approaches --- and remain
unreliable on the evaluative subset that motivates their use in the
first place.

\subsection{Vocabulary Drift Under Few-Shot
Prompting}\label{sec:results-drift}

Two anomalies in Tables~\ref{tab:multilabel}
and~\ref{tab:extractive-evaluative} warrant separate analysis:
GPT-4o-mini's few-shot macro-F1 (0.003) and Mistral's few-shot
macro-F1 (0.016) are approximately two orders of magnitude below the
same models' zero-shot performance.  Ordinary explanations ---
degraded reasoning under more context, distraction by examples --- do
not fit, because the same models retain competitive
acceptance-classification accuracy under few-shot prompting (0.774
and 0.808 respectively; Table~\ref{tab:acceptance}).  Something is
happening specifically to the model's flag output, not its underlying
task understanding.

Manual inspection of the raw prediction output identifies the cause:
under few-shot prompting, both models frequently emit flag names that
are outside the enumerated vocabulary supplied in the prompt.  Rather
than the vocabulary-form \texttt{poor\_code\_quality}, the model
outputs ``Poor code quality'' (case- and separator-inconsistent);
rather than \texttt{manageable\_number\_of\_changed\_lines}, the model
outputs \texttt{managerable\_number\_of\_changed\_lines} (typo) or
\texttt{clear\_commit\_messages} (pluralized).  Occasionally the model
emits an entirely invented flag name not present in the taxonomy at
all.  Because per-flag F1 is computed by exact string match against
the taxonomy, every such emission counts as a false negative for the
correct flag and yields no true positive.
Table~\ref{tab:vocab-drift} quantifies this behavior.

% >>> PASTE TABLE 5.4 HERE  (table_5_4_vocab_drift.tex) <<<

Under few-shot prompting, GPT-4o-mini emits at least one
non-vocabulary flag in \textbf{99.7\% of predictions} (296 of 297
rows), and Mistral does so in \textbf{98.7\% of predictions} (293 of
297).  To distinguish this failure mode from genuine reasoning
failure, we recompute per-flag F1 after applying a conservative
normalization function that maps case- and separator-variants to their
vocabulary form (e.g.,\ ``Poor code
quality''~$\rightarrow$~\texttt{poor\_code\_quality}) but leaves typos
and invented flag names unchanged.  Under normalization,
GPT-4o-mini's few-shot macro-F1 recovers from 0.003 to 0.409, and
Mistral's from 0.016 to 0.402 --- restoring both to a level comparable
with their zero-shot performance.

Two other models exhibit essentially no such drift: \textbf{Claude
Haiku 4.5 emits non-vocabulary flags in at most 0.3\% of predictions
in any condition}, and \textbf{DeepSeek R1 in at most 6.1\% (few-shot)},
with zero-shot DeepSeek R1 drifting on only 0.7\% of predictions.
The gap between drift-prone and drift-free models does not track
model size or open-vs-closed provenance: GPT-4o-mini and Haiku 4.5
are the two smallest closed-source models in our set, yet exhibit
opposite drift behavior.  We interpret this finding as evidence that
instruction-following robustness under longer in-context contexts is
a distinct dimension of model capability from either semantic
reasoning or model scale, and one that materially affects how much
can be inferred about a model's flag-generation capacity from summary
metrics alone.

For the remainder of this paper, we report both the strict and
normalized macro-F1 values for the two drift-affected models where
relevant.  The tables and analyses in
Sections~\ref{sec:results-rq1}--\ref{sec:results-rq3} use the strict
values, matching the standard evaluation protocol; the normalization
above is presented as a diagnostic, not as a substitute for the
strict evaluation.
